% Para utilizar este template siga o tutorial disponível em http://www.biblioteca.ufc.br/wp-content/uploads/2015/09/tutorial-sharelatex.pdf

%%%%%%%%%%%%%%%%%%%%%%%%%%%%%%%%%%%%%%%%%%%%%%%%%%%%%%%
%% Você deve criar uma conta no Overleaf. Depois,    %%
%% vá nas opções no canto esquerdo superior da tela  %%
%% e clique em "Copiar Projeto". Dê um novo nome pa- %%
%% ra o projeto.                                     %%
%%                                                   %%
%% Os principais desenvolvedores deste template são: %%
%%                                                   %%
%%            Ednardo Moreira Rodrigues              %%
%%       (Doutor em Engenharia Elétrica - UFC)       %%
%%                      &                            %%
%%            Alan Batista de Oliveira               %%
%%           (Engenheiro Eletricista - UFC)          %%
%%                                                   %%
%% Revisão:                                          %%
%%                                                   %%
%% - Francisco Edvander Pires Santos;                %%
%% - Juliana Soares Lima;                            %%
%% - Izabel Lima dos Santos;                         %%
%% - Kalline Yasmin Soares Feitosa.                  %%
%% - Eliene Maria Vieira de Moura;                   %%
%%                                                   %%
%% Colaboradores                                     %%
%%                                                   %%
%% -Andrei Bosco Bezerra Torres                      %% 
%% (Professor - Sistemas e Mídias Digitais -         %%
%% Instituto Universidade Virtual - UFC)             %%
%% Tiago ALves Lima                                  %% 
%% (Aluno de Mestrado em Eng. Elétrica)              %%
%%                                                   %%
%% Grande parte do trabalho foi adaptado do template %%
%% da UECE elaborado por:                            %%
%% Thiago Nascimento  (UECE)                         %%
%% Project available on:                             %%
%% https://github.com/thiagodnf/uecetex2             %%
%%                                                   %%
%% "Dúvidas, esclarecimentos ou sugestões podem ser  %%
%% enviadas para o seguinte e-mail:                  %%
%%                                                   %%
%%             atendimentobch@ufc.br                 %%
%%                                                   %%
%% As últimas atualizações estão descritas no inicio %%
%% do arquivo "README.md".                           %%
%%                                                   %%
%%%%%%%%%%%%%%%%%%%%%%%%%%%%%%%%%%%%%%%%%%%%%%%%%%%%%%%

% Importações de pacotes
\usepackage[utf8]{inputenc}                         % Acentuação direta
\usepackage[T1]{fontenc}                            % Codificação da fonte em 8 bits
\usepackage{graphicx}                               % Inserir figuras
\usepackage{amsfonts, amssymb, amsmath}             % Fonte e símbolos matemáticos
\usepackage{booktabs}                               % Comandos para tabelas
\usepackage{verbatim}                               % Texto é interpretado como escrito no documento
\usepackage{multirow, array}                        % Múltiplas linhas e colunas em tabelas
\usepackage{indentfirst}                            % Endenta o primeiro parágrafo de cada seção.
\usepackage{listings}                               % Utilizar codigo fonte no documento
\usepackage{xcolor}
\usepackage{microtype}                              % Para melhorias de justificação?
\usepackage[portuguesekw,ruled,linesnumbered]{algorithm2e}    % Escrever algoritmos
\usepackage{algorithmic}                            % Criar Algoritmos  
%\usepackage{float}                                 % Utilizado para criação de floats
\usepackage{amsgen}
\usepackage{lipsum}                                 % Usar a simulação de texto Lorem Ipsum
%\usepackage{titlesec}                              % Permite alterar os títulos do documento
\usepackage{tocloft}                                % Permite alterar a formatação do Sumário
\usepackage{etoolbox}                               % Usado para alterar a fonte da Section no Sumário
\usepackage[nogroupskip,nonumberlist]{glossaries}   % Permite fazer o glossario

\usepackage[font=singlespacing,format=plain,justification=justified,skip=0pt,singlelinecheck = false]{caption}            % Altera o comportamento da tag caption

\usepackage[alf, abnt-emphasize=bf, recuo=0cm, abnt-etal-cite=2, abnt-etal-list=0, abnt-etal-text=it]{abntex2cite}  % Citações padrão ABNT
%\usepackage[bottom]{footmisc}                      % Mantém as notas de rodapé sempre na mesma posição
%\usepackage{times}                                 % Usa a fonte Times
%%%%%%%%%%%%%%%%%%% AVISO %%%%%%%%%%%%%%%%%%%%%%%%%%%%%%%%%%%%%%%%
%descomente as duas linhas abaixo para alterar o texto de Times New Roman para Arial:

%\usepackage{helvet}
%\renewcommand{\familydefault}{\sfdefault}  % Usa a fonte Arial              
%%%%%%%%%%%%%%%%%%%%%%%%%%%%%%%%%%%%%%%%%%%%%%%%%%%%%%%%%%%%%%%%%%

\usepackage{mathptmx}         % Usa a fonte Times New Roman			%\usepackage{lmodern}         % Usa a fonte Latin Modern
%\usepackage{subfig}          % Posicionamento de figuras
%\usepackage{scalefnt}        % Permite redimensionar tamanho da fonte
%\usepackage{color, colortbl} % Comandos de cores
%\usepackage{lscape}          % Permite páginas em modo "paisagem"
%\usepackage{ae, aecompl}     % Fontes de alta qualidade
%\usepackage{picinpar}        % Dispor imagens em parágrafos
%\usepackage{latexsym}        % Símbolos matemáticos
%\usepackage{upgreek}         % Fonte letras gregas
\usepackage{appendix}         % Gerar o apendice no final do documento
\usepackage{paracol}          % Criar paragrafos sem identacao
\usepackage{lib/ufctex}	      % Biblioteca com as normas da UFC para trabalhos academicos
\usepackage{pdfpages}         % Incluir pdf no documento
\usepackage{amsmath}          % Usar equacoes matematicas
\usepackage{amsthm}
\usepackage{float}
\usepackage{tikz}
\usetikzlibrary{arrows.meta}
\usepackage{pgfplots}
\usepackage[most]{tcolorbox}
\usepackage{mathtools}

\usepackage{caption}
\usepackage{subcaption}

\DeclarePairedDelimiter{\ceil}{\lceil}{\rceil}
\DeclarePairedDelimiter{\floor}{\lfloor}{\rfloor}

\newtheorem{theorem}{Teorema}[chapter]
\newtheorem{corollary}{Corolário}[chapter]
\newtheorem{lemma}{Lema}[chapter]
\newtheorem{proposition}{Proposição}[chapter]
\newtheorem{definition}{Definição}[chapter]
\newtheorem{example}{Exemplo}[chapter]
\newtheorem{pro}{Problema}[chapter]
%AQUI EU FIZ
\newtheorem*{remark}{Conjectura 1}
\newtheorem*{remarkk}{Conjectura 2}

\makeglossaries % Organiza e gera a lista de abreviaturas, simbolos e glossario
\makeindex      % Gera o Indice do documento         

\renewcommand{\labelitemi}{$\bullet$}

\newcommand{\pname}[1]{\textsc{#1}}

\newtcolorbox[auto counter]{ProblemBox}[2][]{
colframe=black!20!white,
%fonttitle=\bfseries,
coltitle=black,
arc=1.5mm,
boxrule=.2mm,
colbacktitle=white,
colback=white,
enhanced,
adjusted title=flush left,
attach boxed title to top center={yshift=-3.5mm,xshift=-4.5cm},
boxed title style={colframe=white,righttitle=-3mm,lefttitle=-3mm},
title=#2,#1
}

\newcommand{\problemDef}[5]{
\vspace{-2mm}
\begin{ProblemBox}[label={#5},nameref={\pname{#1}}]{\pname{#1}}
\begin{tabularx}{\textwidth}{rX}
\textnormal{Entrada:} &\textnormal{#2} \\
            \textnormal{Problema:} & \textnormal{#3}\\
            &\textnormal{#4}
\end{tabularx}
\vspace{-3,6mm}
\end{ProblemBox}
}

\newcommand{\alv}[1]{\textcolor{blue}{(alvaro: #1)}}
\newcommand{\kar}[1]{\textcolor{red}{(karol: #1)}}

\newcommand{\pn}{$P_3$}
\newcommand{\ext}[1]{$Ext(#1)$}
\newcommand{\infd}[1]{$\overrightarrow{I_{P_3}}(#1)$}
\newcommand{\pd}{$\overrightarrow{P_3}$ }
\newcommand{\inum}[1]{$in_{#1}$}
\newcommand{\hnum}[1]{$hn_{#1}$}
\newcommand{\inumd}[1]{$\overrightarrow{in_{#1}}$}
\newcommand{\hnumd}[1]{$\overrightarrow{hn_{#1}}$}

\newcommand{\npct}{$\mathsf{NP}$-completude}
\newcommand{\npc}{$\mathsf{NP}$-completo }
\newcommand{\np}{$\mathsf{NP}$ }
\newcommand{\ind}{\textup{\textsc{Número de Intervalo \inumd{P_3} }}}
\newcommand{\hnd}{\textup{\textsc{Número de Envoltória \hnumd{P_3} }}}
\newcommand{\sam}{\textup{\textsc{Planar SAT-AM3 }}}
